\documentclass[11pt,a4paper]{article}

% Minimale Autorenvorlage fuer das taegliche Schreiben.
% Diese Vorlage schlank halten; fuer die Abgabe in die Publisher-Vorlage ueberfuehren.

\usepackage[utf8]{inputenc}
\usepackage[T1]{fontenc}
\usepackage[ngerman]{babel}
\usepackage{lmodern}
\usepackage{microtype}
\usepackage[a4paper,margin=2.5cm]{geometry}
\usepackage{setspace}
\usepackage{graphicx}
\usepackage{booktabs}
\usepackage{amsmath,amssymb}
\usepackage[numbers,sort&compress]{natbib}
\usepackage[hidelinks]{hyperref}

\onehalfspacing

\title{Arbeitstitel (DE)}
\author{Autorname}
\date{\today}

\begin{document}
\maketitle

\begin{abstract}
Hier die Zusammenfassung einfuegen. Bereits im Entwurf zitieren, z. B. \cite{Weilkiens2024_SysMLv2Book}.
\end{abstract}

\textbf{Schluesselwoerter:} SysML v2, MBSE, Architektur

\section{Einleitung}
Direkt in LaTeX schreiben.

\section{Vorgehen}
Methode und Evidenz beschreiben.

\section{Diskussion}
Ergebnisse und Grenzen einordnen.

\section{Fazit}
Kernaussagen und naechste Schritte zusammenfassen.

\bibliographystyle{plainnat}
\bibliography{../../bibliography/references}

\end{document}
